% =====================================================================
% Appendix A -- Designed but not executed
% =====================================================================
\documentclass[../preamble.tex]{subfiles}
\begin{document}

\chapter{Designed but not executed}
\label{app:designed-not-executed}

The v3 methodology specification describes several pipeline features that
were implemented in the trainer scaffold but were not used for the final
reported results. Each is listed below with its design rationale, the reason
for deferral, and the code location for re-enablement. None of these features
are claimed as results in the body of the thesis.

\section{Exponential Moving Average (EMA)}
\label{app:designed-ema}

The trainer scaffold supports an EMA shadow with $\beta = 0.999$, updated
after each training step. EMA is a low-cost regulariser that often improves
calibration without affecting accuracy. The fine-tune runs that produced the
reported checkpoints did not enable EMA. To re-enable it, set
\texttt{training.ema.enabled=true} and \texttt{training.ema.beta=0.999} in
the configuration. The EMA shadow would need to be swapped in at evaluation
time.

\section{Stochastic Weight Averaging (SWA)}
\label{app:designed-swa}

SWA averages the parameters of the last $N$ checkpoints along the training
trajectory. The intended SWA buffer size was five checkpoints, with the
averaged weights evaluated on the validation set. SWA was not used in the
fine-tune runs because the validation set is small and the gain in
calibration from SWA over temperature scaling was not separately measured.
SWA can be re-enabled via \texttt{training.swa.enabled=true} and
\texttt{training.swa.keep\_last\_n=5}.

\section{$k$-fold cross-validation}
\label{app:designed-kfold}

The methodology specification describes a sweep-then-$k$-fold design: HPO on
the top-$k$ finalists followed by $k$-fold cross-validation to estimate
performance stability. The reported results use a single stratified 80/10/10
split on the Figshare data and a single 2{,}560 / 640 split on PCOSgen, with
a separate held-out test set. Repeated cross-validation would be the right
way to estimate the variance of the ensemble metrics; the thesis instead
relies on the bootstrap 95\% confidence interval on the held-out test set.
Re-enabling $k$-fold requires choosing $k$ (typically 5) and a unified
Figshare-plus-PCOSgen corpus, which is not available without combining the
two datasets under a documented patient-independent rule.

\section{Layer-wise Relevance Propagation (LRP) and SHAP}
\label{app:designed-lrp-shap}

The original FYDP brief listed Grad-CAM, LRP, and SHAP as the XAI
instruments. Only Grad-CAM was executed on the fine-tuned checkpoints. LRP
propagates relevance backward through the network and is well suited to
objective perturbation scoring \cite{samek2017}. SHAP provides pixel-level
attribution with a game-theoretic foundation. Both would extend the
qualitative Grad-CAM analysis in \cref{ch:trust}.

\section{Nine-by-six preprocessing grid}
\label{app:designed-9x6}

The brief described a six-condition preprocessing grid (CLAHE on/off,
Gaussian $\times$ three $\sigma$, SRAD on/off). The executed ablation
reduced this to two conditions (SRAD vs.\ Gaussian) with CLAHE fixed. The
choice was deliberate: isolating a single variable made it possible to
attribute any performance change to the denoiser rather than to a
combination of contrast and noise changes. The full nine-by-six grid is
available for re-enablement in the trainer configuration.

\section{InceptionV3}
\label{app:designed-inception}

The brief listed InceptionV3 as one of the nine architectures. The executed
roster substitutes Swin-Tiny and MobileNetV3-Large, which were the
contemporary transformer and mobile architectures available in
\texttt{timm}. InceptionV3 was removed to keep the parameter-count envelope
comparable across the roster.

\end{document}